
Il suffit donc, pour achever la démonstration, de noter que :
\[
\Supp\bigl(H^*(\underline{M}^\bullet\overset{\mathbb L}{\otimes}\underline{M}'{}^\bullet)\bigr)
 \subset \Supp(H^*(\underline{M}^\bullet))\cap\Supp(H^*(\underline{M}'{}^\bullet)).
\]

Pour montrer que $K^{\bullet}(X)_{d+1}$ est nul, on prend $x\in K^{\bullet}(X)_{d+1}$, et $Y=X$ il existe alors $\underline{M}^\bullet$ tel que $x=\cl(\underline{M}^\bullet)$, et $\codim_X(\Supp(H^*(\underline{M}^\bullet)))\ge d+1$; comme $X$ est de dimension $d$, cela entraine $\Supp(H^*(\underline{M}^\bullet))=\varnothing$, donc $\underline{M}^\bullet$ est acyclique, et $x=0$.

\textbf{Lemme 6.7.} \uline{Soit} $X=\Spec(A)$ \uline{un schéma affine noethérien}, $x_1,\ldots,x_k$ \uline{un nombre fini de points de} $X$, $\underline{E}$ \uline{un} $\underline{O}_X$\uline{-Module inversible}. \uline{Alors}, \uline{il existe une section} $s\in\Gamma(X,\underline{E})$ \uline{telle que pour tout} $i$, $s_{x_i}$ \uline{soit une base de} $\underline{E}_{x_i}$.

Posons $M=\Gamma(X,\underline{E})$. Pour tout $i$, soit $x'_i$ un point fermé de $X$ tel que $x'_i\in\overline{\{x_i\}}$. Les $x'_i$ correspondent à des idéaux maximaux $\underline{m}_i$ de $A$, qu'on peut supposer tous distincts, quitte à changer l'ensemble d'indices de façon que chacun des $\underline{m}_i$ intervienne une fois et une seule dans la famille $(\underline{m}_i)$. Comme $\underline{E}$ est inversible, on a pour tout $i$ un isomorphisme $M_{\underline{m}_i}\simeq A_{\underline{m}_i}$ d'où on déduit un isomorphisme $M/\underline{m}_iM\simeq A/\underline{m}_i$, puis $\prod_i M/\underline{m}_iM\simeq\prod_i A/\underline{m}_i$, soit encore $A/\prod_i\underline{m}_i\simeq M/(\prod_i\underline{m}_i)M$. En relevant ce dernier isomorphisme, on obtient un homomorphisme $A\to M$ tel que l'homomorphisme $A_{\underline{m}_i}\to M_{\underline{m}_i}$ correspondant soit pour tout $i$ un isomorphisme modulo $\underline{m}_i$, donc un isomorphisme d'après Nakayama. L'image de $1$ par cet homomorphisme donne donc une section $s$ de $\underline{E}$ qui est une base de $\underline{E}_{x'_i}$ pour tout $i$. Comme l'ensemble des points où $s$ est une base de $\underline{E}$ est un ouvert de $X$, $s$ est aussi une base de $\underline{E}_{x_i}$ pour $i=1,\ldots,k$.

\textbf{Lemme 6.8.} \uline{Soient} $X$ \uline{un schéma noethérien ayant un faisceau ample} $\underline{L}$, $\underline{E}$ \uline{un} $\underline{O}_X$\uline{-Module inversible}, $x_1\ldots,x_k$ \uline{des points de} $X$. \uline{Alors il existe un}
